\documentclass[10pt,a4paper]{amsart}
\usepackage{amsmath,amssymb,amsthm,graphicx}
% Retained minimal preamble for the Stacks Project chapter flat.tex (More on Flatness).
% Verbatim excerpt of the chapter preamble: Xy-pic setup (native preamble lines 14-18) and the native
% theorem declarations with the mathematical macros (native preamble lines 154-217). Copied mechanically;
% nothing added, removed or redefined. The wrapper class and the amsmath/amssymb/amsthm/graphicx base
% packages are supplied by the assembler, so the class/externaldocument part of the native preamble is omitted.
% For commutative diagrams we use Xy-pic
\usepackage[all]{xy}

% We use 2cell for 2-commutative diagrams.
\xyoption{2cell}
% Theorem environments.
%
\theoremstyle{plain}
\newtheorem{theorem}[subsection]{Theorem}
\newtheorem{proposition}[subsection]{Proposition}
\newtheorem{lemma}[subsection]{Lemma}

\theoremstyle{definition}
\newtheorem{definition}[subsection]{Definition}
\newtheorem{example}[subsection]{Example}
\newtheorem{exercise}[subsection]{Exercise}
\newtheorem{situation}[subsection]{Situation}

\theoremstyle{remark}
\newtheorem{remark}[subsection]{Remark}
\newtheorem{remarks}[subsection]{Remarks}

\numberwithin{equation}{subsection}

% Macros
%
\def\lim{\mathop{\mathrm{lim}}\nolimits}
\def\colim{\mathop{\mathrm{colim}}\nolimits}
\def\Spec{\mathop{\mathrm{Spec}}}
\def\Hom{\mathop{\mathrm{Hom}}\nolimits}
\def\Ext{\mathop{\mathrm{Ext}}\nolimits}
\def\SheafHom{\mathop{\mathcal{H}\!\mathit{om}}\nolimits}
\def\SheafExt{\mathop{\mathcal{E}\!\mathit{xt}}\nolimits}
\def\Sch{\mathit{Sch}}
\def\Mor{\mathop{\mathrm{Mor}}\nolimits}
\def\Ob{\mathop{\mathrm{Ob}}\nolimits}
\def\Sh{\mathop{\mathit{Sh}}\nolimits}
\def\NL{\mathop{N\!L}\nolimits}
\def\CH{\mathop{\mathrm{CH}}\nolimits}
\def\proetale{{pro\text{-}\acute{e}tale}}
\def\etale{{\acute{e}tale}}
\def\QCoh{\mathit{QCoh}}
\def\Ker{\mathop{\mathrm{Ker}}}
\def\Im{\mathop{\mathrm{Im}}}
\def\Coker{\mathop{\mathrm{Coker}}}
\def\Coim{\mathop{\mathrm{Coim}}}

% Boxtimes
%
\DeclareMathSymbol{\boxtimes}{\mathbin}{AMSa}{"02}

%
% Macros for moduli stacks/spaces
%
\def\QCohstack{\mathcal{QC}\!\mathit{oh}}
\def\Cohstack{\mathcal{C}\!\mathit{oh}}
\def\Spacesstack{\mathcal{S}\!\mathit{paces}}
\def\Quotfunctor{\mathrm{Quot}}
\def\Hilbfunctor{\mathrm{Hilb}}
\def\Curvesstack{\mathcal{C}\!\mathit{urves}}
\def\Polarizedstack{\mathcal{P}\!\mathit{olarized}}
\def\Complexesstack{\mathcal{C}\!\mathit{omplexes}}
% \Pic is the operator that assigns to X its picard group, usage \Pic(X)
% \Picardstack_{X/B} denotes the Picard stack of X over B
% \Picardfunctor_{X/B} denotes the Picard functor of X over B
\def\Pic{\mathop{\mathrm{Pic}}\nolimits}
\def\Picardstack{\mathcal{P}\!\mathit{ic}}
\def\Picardfunctor{\mathrm{Pic}}
\def\Deformationcategory{\mathcal{D}\!\mathit{ef}}

% --- appended 2026-10-12 (Chief): the project own definitions of the reference/slogan/history
% display environments, copied mechanically from the shared Stacks preamble lines 8-12
% (\usepackage{verbatim} plus the three \newenvironment lines, which the chapter preamble keeps
% verbatim from the project preamble). Nothing added, removed or redefined: these environments
% carry no mathematical content and the project itself typesets them as comments. Without these
% three lines a counted statement whose own text uses them fails with Environment slogan undefined.
\usepackage{verbatim}
\newenvironment{reference}{\comment}{\endcomment}
\newenvironment{slogan}{\comment}{\endcomment}
\newenvironment{history}{\comment}{\endcomment}
\usepackage{cleveref}
\crefname{originallemma}{original lemma}{original lemmas}
\Crefname{originallemma}{Original lemma}{Original lemmas}
\makeatletter
\newlabel{lemma-affine-locally-projective-pure}{{05K6}{source}{}{sourcepointer-1}{}}
\newlabel{lemma-affine-locally-projective-pure@cref}{{[originallemma][0][]05K6}{[1][0][]source}{}{}{}}
\makeatother
\crefname{originallemma}{original lemma}{original lemmas}
\Crefname{originallemma}{Original lemma}{Original lemmas}
\makeatletter
\newlabel{lemma-criterion}{{05L4}{source}{}{sourcepointer-2}{}}
\newlabel{lemma-criterion@cref}{{[originallemma][0][]05L4}{[1][0][]source}{}{}{}}
\makeatother
\crefname{originallemma}{original lemma}{original lemmas}
\Crefname{originallemma}{Original lemma}{Original lemmas}
\makeatletter
\newlabel{lemma-finite-presentation-flat-along-fibre}{{05KW}{source}{}{sourcepointer-3}{}}
\newlabel{lemma-finite-presentation-flat-along-fibre@cref}{{[originallemma][0][]05KW}{[1][0][]source}{}{}{}}
\makeatother
\crefname{originallemma}{original lemma}{original lemmas}
\Crefname{originallemma}{Original lemma}{Original lemmas}
\makeatletter
\newlabel{more-morphisms-lemma-relative-assassin-constructible}{{05KR}{source}{}{sourcepointer-4}{}}
\newlabel{more-morphisms-lemma-relative-assassin-constructible@cref}{{[originallemma][0][]05KR}{[1][0][]source}{}{}{}}
\makeatother
\crefname{originallemma}{original lemma}{original lemmas}
\Crefname{originallemma}{Original lemma}{Original lemmas}
\makeatletter
\newlabel{morphisms-lemma-constructible-containing-open}{{05LW}{source}{}{sourcepointer-5}{}}
\newlabel{morphisms-lemma-constructible-containing-open@cref}{{[originallemma][0][]05LW}{[1][0][]source}{}{}{}}
\makeatother
\title{The locus of points over which a flat finitely presented module is pure along the fibre is open}
\author{The Stacks Project Authors}
\date{stacks-project:flat.tex:local-snapshot-20261011; GNU FDL 1.2 or later}
\begin{document}
\maketitle
\noindent\textit{Editorial note.} All delivered mathematical text is the author's own, copied from the pinned Stacks Project chapter file flat.tex; no statement or proof step was rewritten, reordered or removed.\par\medskip
\noindent\textit{Editorial note.} Preamble provenance: the chapter's own retained minimal preamble (sources/stacks-flat/retained-preamble.tex, lines 1-85) is a verbatim excerpt of the native Xy-pic setup and of the native theorem declarations and mathematical macros, followed by the project's own reference/slogan/history display-environment definitions copied mechanically from the shared preamble; the wrapper class and the amsmath/amssymb/amsthm/graphicx base packages are supplied by the assembler, and extra\_packages is left empty because everything needed is already in the retained preamble. No macro was invented or redefined.\par\medskip
\noindent\textit{Editorial note.} Same-chapter pointers are located at their \textbackslash{}label lines inside the pinned chapter file; cross-chapter pointers are located in the project's own tag list staged in the snapshot (sources/stacks-chow/tags.txt), at the verbatim lines carrying their canonical tags.\par\medskip
\noindent\textit{Allowed standard prerequisites (editorial).}\begin{itemize}
\item The affine etale neighbourhood with projective global sections. Same chapter flat.tex, label lemma-finite-presentation-flat-along-fibre, invoked at proof line 5502 and claimed separately as S14-algebra-fp-flat-along-fibre. Produces the elementary etale neighbourhood and the affine etale neighbourhood whose global sections of the pulled back module are projective over the ring of the base.
\item Purity of a module with projective global sections, the purity criterion, and constructibility of the relative assassin condition. Same chapter flat.tex, labels lemma-affine-locally-projective-pure (proof line 5515) and lemma-criterion (proof lines 5526 and 5532); More on Morphisms, label more-morphisms-lemma-relative-assassin-constructible (proof line 5524). The first makes the pulled back module universally pure, the constructibility statement makes the set of points where the associated points of the fibre lie in the image of the etale morphism constructible, and the criterion converts that inclusion into purity, which is how the open neighbourhood is produced.
\item A constructible set containing the spectrum of the local ring at a point contains an open neighbourhood, and the image of an open subset under an etale morphism is open. Morphisms, label morphisms-lemma-constructible-containing-open (proof line 5529). Turns the containment of the local spectrum in the constructible set into an open neighbourhood of the point, whose image in the base is the desired open set of points of purity.
\end{itemize}\medskip
\ifdefined\hypertarget\hypertarget{sourcepointer-1}{}\fi
\noindent\textit{Source reference.} Original lemma 05K6: sources/stacks-flat/pinned-original-source.tex, source lines 5148--5148. Navigation pointer for \textbackslash{}ref\{lemma-affine-locally-projective-pure\} used inside the counted span (line 5515); the argument that consumes it is retained in the span.\par\medskip
\ifdefined\hypertarget\hypertarget{sourcepointer-2}{}\fi
\noindent\textit{Source reference.} Original lemma 05L4: sources/stacks-flat/pinned-original-source.tex, source lines 5197--5197. Navigation pointer for \textbackslash{}ref\{lemma-criterion\} used inside the counted span (line 5526); the argument that consumes it is retained in the span.\par\medskip
\ifdefined\hypertarget\hypertarget{sourcepointer-3}{}\fi
\noindent\textit{Source reference.} Original lemma 05KW: sources/stacks-flat/pinned-original-source.tex, source lines 3342--3342. Navigation pointer for \textbackslash{}ref\{lemma-finite-presentation-flat-along-fibre\} used inside the counted span (line 5502); the argument that consumes it is retained in the span.\par\medskip
\ifdefined\hypertarget\hypertarget{sourcepointer-4}{}\fi
\noindent\textit{Source reference.} Original lemma 05KR: sources/stacks-chow/tags.txt, source lines 7210--7210. Navigation pointer for \textbackslash{}ref\{more-morphisms-lemma-relative-assassin-constructible\} used inside the counted span (line 5524); the argument that consumes it is retained in the span.\par\medskip
\ifdefined\hypertarget\hypertarget{sourcepointer-5}{}\fi
\noindent\textit{Source reference.} Original lemma 05LW: sources/stacks-chow/tags.txt, source lines 7250--7250. Navigation pointer for \textbackslash{}ref\{morphisms-lemma-constructible-containing-open\} used inside the counted span (line 5529); the argument that consumes it is retained in the span.\par\medskip
\begin{lemma}
\label{lemma-flat-finite-presentation-purity-open}
Let $f : X \to S$ be a morphism of finite presentation.
Let $\mathcal{F}$ be a quasi-coherent $\mathcal{O}_X$-module
of finite presentation flat over $S$. Then the set
$$
U = \{s \in S \mid \mathcal{F}\text{ is pure along }X_s\}
$$
is open in $S$.
\end{lemma}

\begin{proof}
Let $s \in U$. Using
Lemma \ref{lemma-finite-presentation-flat-along-fibre}
we can find an elementary \'etale neighbourhood
$(S', s') \to (S, s)$ and a commutative diagram
$$
\xymatrix{
X \ar[d] & X' \ar[l]^g \ar[d] \\
S & S' \ar[l]
}
$$
such that $g$ is \'etale, $X_s \subset g(X')$, the schemes
$X'$, $S'$ are affine, and such that $\Gamma(X', g^*\mathcal{F})$
is a projective $\Gamma(S', \mathcal{O}_{S'})$-module.
Note that $g^*\mathcal{F}$ is universally pure over $S'$, see
Lemma \ref{lemma-affine-locally-projective-pure}.
Set $W' \subset X \times_S S'$ equal to the image of the \'etale morphism
$X' \to X \times_S S'$. Note that $W$ is open and quasi-compact over
$S'$. Set
$$
E = \{t \in S' \mid \text{Ass}_{X_t}(\mathcal{F}_t) \subset W' \}.
$$
By
More on Morphisms,
Lemma \ref{more-morphisms-lemma-relative-assassin-constructible}
$E$ is a constructible subset of $S'$. By
Lemma \ref{lemma-criterion}
we see that $\Spec(\mathcal{O}_{S', s'}) \subset E$.
By
Morphisms, Lemma \ref{morphisms-lemma-constructible-containing-open}
we see that $E$ contains an open neighbourhood $V'$ of $s'$.
Applying
Lemma \ref{lemma-criterion}
once more we see that for any point $s_1$ in the image of $V'$ in $S$
the sheaf $\mathcal{F}$ is pure along $X_{s_1}$. Since
$S' \to S$ is \'etale the image of $V'$ in $S$ is open and we win.
\end{proof}

\end{document}
